\documentclass[11pt]{article}
\usepackage[utf8]{inputenc}
\usepackage[T1]{fontenc}
\usepackage{lmodern}
\usepackage{amsmath,amssymb}
\usepackage[margin=1in]{geometry}
\usepackage{booktabs,tabularx,xltabular,array,enumitem,xcolor}
\usepackage[colorlinks=true,linkcolor=blue!50!black,urlcolor=blue!50!black]{hyperref}
\setlist{itemsep=2pt,topsep=4pt}
\newcolumntype{L}{>{\raggedright\arraybackslash}X}
\title{How we went about this\\[0.5ex]\large The planning session: a process report on planning a textbook for one research paper}
\author{Prepared by Claude, Claude Code session\\ on branch \texttt{claude/textbook-paper-prerequisites-clhh5a}}
\date{11 to 15 October 2026; report written 25 October 2026}
\begin{document}
\maketitle

\begin{abstract}
The goal was a textbook that takes a reader with first courses in analysis, linear algebra, complex analysis,
statistics and Riley-Hobson-Bence to the point of reading the note \emph{Local factors of the adjoint
$L$-function and the $\ell$-adic valuations of the modular degree}, and that experts can open as a reference.
Over five days one orchestrating session, working with the author and about forty sub-agents, produced a master
plan, verified outlines of all 22 chapters and an appendix of conventions, a verified bibliography, a record of
80 findings about the paper, a memo for the author, a revised version of the paper with stronger results, and
the tools and instructions to write the full text in a new session. No chapter text has been written yet. This
report describes how the work was organised, what each party did, what went wrong and how it was caught, what it
cost, and what we would do differently.

\emph{Scope of this report.} It covers only the planning session, up to the hand-over of 15 October 2026. After the hand-over, another session continued from \texttt{CONTINUE.md}: it wrote the book and finalised the paper, and its own report is \texttt{final/process/how\_we\_went\_about\_this.pdf} on the same branch. Statements below such as ``no chapter text has been written yet'' describe the state at the hand-over.
\end{abstract}

\section{At a glance}

{\small
\begin{xltabular}{\linewidth}{>{\bfseries}p{3.6cm}L}
\toprule
Goal & A modular textbook (about 22 chapters) from undergraduate background to a specific research paper, also
usable as a reference by experts.\\
Starting material & The paper's LaTeX source (547 lines) and an empty git repository.\\
Approach & Plan centrally; outline every chapter with a separate agent, chapter by chapter at the author's pace;
verify claims computationally and by independent re-derivation; keep a single running record of findings;
write the full text later from the outlines.\\
Agents & One orchestrator; about 40 sub-agents (23 outline agents, 4 housekeeping agents, 1 paper-revision agent,
2 setup agents, and 22 line-up agents of a first attempt that was stopped by a usage limit).\\
Human role & Set the goal and audience, chose the pace and order (usually one to three chapters at a time), took
four binding decisions, and asked for the memo, the revised paper and this report.\\
Verification & PARI/GP (via cypari2) on Cremona's tables to conductor 10{,}000 (and 30{,}000 for one question);
the orchestrator re-derived every high-impact claim before recording it; the published DFG04 was read to settle
the central open question.\\
Main outputs & \texttt{PLAN.md}; 23 outlines (24{,}421 lines); \texttt{docs/book\_plan.pdf} (649 pages);
\texttt{refs.bib} and chapter bib files (203 entries); \texttt{important/important\_notes.md} (80 points, 27
book-internal notes); \texttt{important/paper\_memo.md}; \texttt{paper/note\_v2.tex} (32 pages);
\texttt{CONTINUE.md}, setup script and writing workflow.\\
Not done & Chapter text (all 22 chapters), Appendix B text, the rest of Appendix A, the integration pass.\\
Main open risks & The new proofs in the revised paper need the author's review; one step of the main
strengthening (DFG04's integral structure versus $H^1(X_0(N),\mathbb Z_\ell)$) is sketched, not proved; the
paper's data tables in version 2 are placeholders.\\
\bottomrule
\end{xltabular}
}

\section{Timeline}

{\small
\begin{xltabular}{\linewidth}{p{2.4cm}L}
\toprule
Date & What happened\\
\midrule
11 Oct & Read the paper in full. Wrote the master plan (audience, 22 chapters in six Parts, scope and paper hooks
of every chapter, dependency graph, paper-to-chapter map, writing conventions). Built the LaTeX skeleton, a
lint script and a per-chapter build script. Launched six parallel workflows (one per Part) to outline, write,
review and fix every chapter. Verified the bibliography seed online.\\
11 to 12 Oct & All 22 line-up agents of the first attempt stopped at a session usage limit, having used about
2.2 million tokens without finishing. Five partial outlines and a partial appendix survived. The model was
switched and the author asked for a cost estimate, then chose to continue chapter by chapter.\\
12 to 14 Oct & Completed the five partial outlines, then outlined the remaining chapters in batches of one to
three agents running in parallel. Each finished outline was checked, its findings about the paper verified and
recorded, and committed. A sub-agent installed PARI/GP on 13 October, after which agents checked claims on
Cremona's tables.\\
14 Oct & Appendix B fixed book-wide conventions and listed every inconsistency between outlines. The Chapter 20
agent read the published DFG04 and settled the checklist behind the central strengthening. In parallel: a memo
for the author, convention fixes across the outlines, a revision of the Chapter 21 outline, and bibliography
consolidation.\\
15 Oct & The author took decisions D1 to D4. A revised paper (version 2) was written and compiled. The plan
and outlines were assembled into one PDF. Instructions, a setup script and a writing workflow were added for a
new session.\\
\bottomrule
\end{xltabular}
}

\section{How the work was organised}

\subsection{Planning first, centrally}
The orchestrator read the paper in full before planning anything, then wrote \texttt{PLAN.md}: the audience and
two reading modes (linear, and refresher through a ``Core results at a glance'' card at the head of each chapter);
the design principles (one topic per chapter, every chapter linked to the exact places the paper uses it, three
declared levels of rigour); the scope of each of the 22 chapters, written in enough detail that an agent could
work from it alone; a dependency graph; a table mapping every notion of the paper to chapters; and strict
conventions (file structure, labels, citations, typography). The plan was the single source of truth that every
later agent was told to read first.

\subsection{Infrastructure before content}
A shared LaTeX preamble holds every macro, so chapters define none. A lint script rejects forbidden dashes,
non-ASCII characters, undefined label prefixes and stray commands. A build script compiles one chapter alone
with all other chapters' labels pre-registered, so cross-references resolve. A bibliography agent verified every
seed entry online. This let any agent check its own output mechanically.

\subsection{A first attempt that failed, and what replaced it}
The first attempt ran six workflows at once, one per Part, each pipelining line-up, writing, two reviews and up to
three fix rounds per chapter. All 22 line-up agents hit a session usage limit before finishing, after about 2.2
million tokens. The lesson was that the full pipeline for 22 chapters at once is too large to run unsupervised in
one sitting, and that the outline stage alone was already substantial.

The replacement was incremental and human-paced: the author named the next chapter or batch (``finish chapter 1
then stop'', ``do 7, 8, 9''), the orchestrator launched one agent per chapter, and each batch was checked and
committed before the next. Outlines became the main deliverable of this phase; writing the text was deferred.

\subsection{What each outline agent was asked to do}
Every outline agent received the same brief, adapted to its chapter: read the plan, the paper and every outline
already written; supply exactly what earlier outlines defer to this chapter and nothing they already cover; line
up every result in teaching order with a precise statement, a rigour level, a reference and the paper location
that uses it; give a section plan, a coverage table, exercises and notes for the writer; verify computations by
hand or with Python and later PARI; and end with a section of \emph{new} problems found in the paper, with
location, fix and confidence. Agents running at the same time were told about each other, and the orchestrator
relayed coordination messages between running agents (for example, which chapter owns a lemma).

\subsection{The running record of findings}
At the author's request, a single file, \texttt{important/important\_notes.md}, collected every point about the
paper as it was found, with source, argument and confidence. Book-internal inconsistencies were kept separately
(B1 to B27). Before recording a high-impact claim the orchestrator re-derived it independently (for example the
determinant argument bounding $h_\ell(E)$, the cyclic 27-isogeny argument, the equality coming from the projection
formula, and every curve example by PARI). Corrections were made in place with a note saying what changed and
why, so the file stayed an honest history.

\subsection{Consolidation, then the paper}
When all outlines existed, four agents ran in parallel with disjoint file ownership: a memo for the author
reconciling the 80 points into their final state; application of the conventions of Appendix B to all outlines
and to the plan; a revision of the Chapter 21 outline, which had been written first and was out of date; and a
consolidation of the bibliography so that the full build and every standalone chapter build work. After the
author's decisions, one agent wrote a revised paper incorporating the strengthenings with proofs and every
correction, with a change log of what the author must check.

\subsection{Handover}
The plan and outlines were assembled into one PDF. \texttt{CONTINUE.md} gives the state of every component, the
prompt to start a new conversation, the setup steps, the pilot, the writing order and the rules. A setup script
installs the toolchain and data in a fresh container, and a workflow script runs the write, review and fix loop
from the existing outlines.

\section{Who did what}

{\small
\begin{xltabular}{\linewidth}{>{\bfseries}p{2.6cm}L}
\toprule
The author & Set the goal, the audience and the use as a reference; decided pace and order; asked to keep the
important points in a separate file; took the binding decisions D1 (state the book's stronger results as
theorems), D2 (length is a guideline, not a cap, since the book is for online reading), D3 (the better character
convention, $\chi_{q^*}$, with the paper to be rewritten accordingly) and D4 (Manin constant positive); asked for
the cost estimate, the memo, the revised paper, the plan PDF and this report.\\
Orchestrator & Read the paper; wrote the plan, tools and instructions; wrote each agent's brief; checked every
returned file (length, typography, structure); re-derived high-impact claims before recording them; relayed
coordination between agents; fixed its own planning errors when agents found them; committed and pushed after
every step; reported to the author and asked before steps the author wanted to control.\\
Outline agents & Read the plan, the paper and all earlier outlines; produced the outlines; computed and checked
examples; reported new findings and inconsistencies.\\
Other agents & Bibliography verification and consolidation; convention fixes; Chapter 21 revision; the memo; the
revised paper; the preface and partial Appendix A.\\
\bottomrule
\end{xltabular}
}

\section{What went wrong, and how it was caught}

\begin{itemize}
\item \textbf{Usage limit in the first attempt.} All parallel line-up agents stopped at once. Partial files were
committed as work in progress and completed one by one.
\item \textbf{Errors in the orchestrator's own plan and notes,} each caught by a later agent and corrected in place:
the Manin constant on the wrong side of the Zagier-Watkins formula; the wrong zeta factor in the naive adjoint
$L$-function; a wrong reason (``the Hecke algebra is reduced'') in note 19, since $U_p$ can have Jordan blocks;
the wrong Kodaira type for a twist in note 15; the wrong linear factor of the $\mathrm I_0^*$ cubic in note 69;
counts quoted for conductor below 10{,}000 that were for 5{,}000; and note 7, which overstated a problem that is
only an imprecision.
\item \textbf{Inconsistency between independently written outlines,} for example two incompatible Frobenius
conventions, the same symbol for two objects, and the same theorem proved in two chapters. Handled by Appendix B,
which fixed sixteen conventions and listed 48 disagreements, and a dedicated fixing pass.
\item \textbf{Outlines written early became stale} as later outlines settled questions. Chapter 21, the first
written, needed a full revision at the end.
\item \textbf{A mechanical fix introduced new errors:} a regular-expression substitution glued two macros together
(\verb|\ad^0M| became \verb|\adM|). Caught by compiling the plan PDF and fixed.
\item \textbf{Minor rule slips by agents:} three agents ran a read-only \texttt{git status} despite being told not
to use git; each reported it.
\item \textbf{Two agents worked from different snapshots,} so one reported stale claims about another's files. The
orchestrator checked the files before acting.
\end{itemize}

\section{What the outlining found in the paper}
Outlining forced every statement of the paper to be restated precisely and checked against its sources and the
data. The result, reconciled in the memo, is in five groups: 12 strengthenings (among them the equality in the main
identity without a multiplicity-one hypothesis, confirmed against the published DFG04; the equality in the
twisting theorem; the removal of an exception at the prime 2; trivial Tamagawa factors at potentially good primes,
answering half of an open question), 10 errors (for example a false twist-invariance claim, a false statement about
Euler factors at 2, a paraphrase of Calegari-Emerton that fails for CM curves, two factual slips about specific
curves), 14 gaps in proofs, 17 expository fixes, and 9 items still open. Version 2 of the paper incorporates them.

\section{Effort and cost indicators}
The figures below are from the sub-agents' own usage reports and the git history; they are indicative, not an
accounting.
\begin{center}\small
\begin{tabular}{lr}
\toprule
Elapsed time, first commit to handover & 11 to 15 October (about 4.5 days, with pauses)\\
Commits & 46\\
Sub-agent tokens per outline & about 240{,}000 to 620{,}000\\
Sub-agent tokens, all outlines & roughly 10 million\\
Sub-agent tokens, first attempt (unfinished) & about 2.2 million\\
Typical outline agent run time & 15 to 60 minutes\\
Outline size & 786 to 1{,}914 lines each; 24{,}421 in total\\
Planned book size & about 49{,}000 lines of LaTeX; roughly 1{,}100 to 1{,}400 pages\\
\bottomrule
\end{tabular}
\end{center}
The orchestrator's own token use is not included. An outline turned out to hold most of a chapter's content:
about half its planned length in lines and about three quarters in words.

\section{Lessons}
\begin{enumerate}
\item \textbf{Plan centrally and in detail before delegating.} A precise scope per chapter, a paper-to-chapter map
and binding conventions let independent agents produce compatible work.
\item \textbf{Do not launch everything at once.} A staged, human-paced sequence with a commit after each step
survived interruptions; the all-at-once attempt did not.
\item \textbf{Make later agents read earlier output.} Each outline agent read every previous outline, which kept
the split between chapters clean and let later agents correct earlier ones.
\item \textbf{Keep one running record of findings, and correct it in place.} It turned scattered observations into
a coherent memo and a revised paper.
\item \textbf{Verify before recording.} Independent re-derivation and computation caught several confident but wrong
claims, including some of the orchestrator's own.
\item \textbf{Fix conventions early.} The Frobenius and notation conflicts would have been cheaper to prevent with a
notation appendix written before the outlines rather than after.
\item \textbf{Give the computational tools to agents from the start.} PARI arrived halfway; the outlines written
after it are better checked.
\item \textbf{Write the hand-over as you go.} The instructions, setup script and workflow make it possible to
continue in a fresh session without this conversation.
\end{enumerate}

\section{State at hand-over and next steps}
Everything is on the branch \texttt{claude/textbook-paper-prerequisites-clhh5a}. The next steps, in
\texttt{CONTINUE.md}: run the setup script; do three small tasks (convert the remaining $\chi_q$ to $\chi_{q^*}$,
update one hedge in the Chapter 18 outline, finish Appendix A); write one pilot chapter end to end with two reviews
and fix rounds, and measure its length, quality and cost; then write the remaining chapters Part by Part, Chapters
21 and 22 last; then the back matter and the integration pass. In parallel, the author reviews the revised paper.

\appendix
\section{Files}
{\small
\begin{xltabular}{\linewidth}{>{\ttfamily}p{5.2cm}L}
\toprule
PLAN.md & Master plan, conventions, binding decisions (Section 7.5).\\
outlines/chNN\_*.md & Outlines of Chapters 1 to 22.\\
outlines/chB\_notation.md & Book-wide conventions and symbol index.\\
docs/book\_plan.pdf & Overview, plan and all outlines in one PDF.\\
important/important\_notes.md & Full record of findings about the paper and the book.\\
important/paper\_memo.md & Reconciled memo for the author.\\
paper/note\_v2.tex, docs/note\_v2.pdf & Revised paper; change log in paper/CHANGES\_v2.md.\\
refs.bib, bib/chNN.bib & Bibliography (shared and single-chapter).\\
CONTINUE.md & How to continue in a new session.\\
tools/ & Build, lint, setup, plan-PDF and chapter-writing workflow scripts.\\
\bottomrule
\end{xltabular}
}
\end{document}
